\section{Constuction of Boneh-Franklin IBE with Chosen Ciphertext Security}

In his original paper \cite{Shamir41} Shamir introduced a public key encryption
scheme that resembled an ideal mail system, in which the public key can be an
arbitrary string which  uniquely identifies the user in a way that he/she
cannot later deny and that is readily available to the other party.\cite{Shamir41}

Shamir's original motivation for identity-based encryption was to simplify the
certificate, signing, and key management processes with regards
for example in a mail system.\cite{BF01}.

%% ttt
Let's say when \textit{Alice} wants to send  a mail to \textit{Bob} at \url{bob@company.com}
she simply encrypts her message using the public key string \lstinline{bob@company.com}. 
In the proposed cryptosystem there is no need for \textit{Alice} to obtain \textit{Bob}’s public key
certificate. When Bob getting noticed that he received the encrypted mail he contacts a trusted
third party called the \textit{Private Key Generator} (denoted as \textit{PKG}  furthermore) and 
after authenticating himself to the \textit{PKG} \textit{Bob} obtains his private key
from the \textit{PKG}. \cite{BF01}

The authentication takes places in the same way \textit{Bob} would authenticate
himself to a trusted Certificate Authority (denoted as CA furthermore).
\textit{Bob} can then therefore read his e-mail sent by \textit{Alice}.

In comparison to an existing secure e-mail infrastructure, with the use of the proposed IBE scheme\cite{BF01}
\textit{Alice} can send an encrypted mail to \textit{Bob} even if he has not yet
setup his public key certificate.


%% todo: ide a kesobbi varians + nehany hatrany, Revocation, Key escrow, ID is public
Also note that key escrow is inherent in identity-based
e-mail systems: the PKG knows Bob’s private key. We discuss key revocation, as well as several new
applications for IBE schemes in the next section.

%% Review previous schemes + new variants
\subsection{Application and variants of Identity-Based Encryption }

% Fujisaki-Okamoto [16]
\subsection{Building blocks for Bonek-Franklin IBE Scheme}


% --------------------------------------